\question
Над булевой функцией разрешены два «зеркала». Зеркало $N$ заменяет результат на противоположный: $g(x)=\neg f(x)$. Зеркало $Z$ заменяет все аргументы на противоположные: $g(x)=f(\overline{x})$. Зеркала можно применять сколько угодно раз и в любом порядке. Пусть $f(x,y,z)=1$ тогда и только тогда, когда хотя бы две из трёх переменных равны $1$. Сколько различных функций можно получить (включая саму $f$) и сколько из них образуют полную систему (каждая по отдельности)? Ответы обоснуйте.

---------------

Автор -- Дарья Гаркуша